% !TEX root = ../main.tex
\clearpage
\section{行业问题与用户需求}
\subsection*{低空产业与金融研究对象}\label{page:05}
\ifdraftversion{\noindent\small\color{gray}写作页位 5\quad 负责人 武桐宇\par}\fi

\textbf{本页结论：}低空研究不能把整机企业单独视作研究对象；材料、动力、控制、整机、基础设施与运营共同决定事件可能落到哪一类公司，因而需要以可追溯的产业关系组织证据。四部门发布的《通用航空装备创新应用实施方案（2024--2030年）》为通用航空装备发展提供了国家层面的政策背景\cite{wty-miit-plan}；广东行动方案则把起降、通信导航监视、数据治理与应用场景并列为建设事项\cite{wty-gd-plan}。

研究对象不是“所有低空概念股”，而是当某一公开事件出现时，能够说明其\emph{事件来源—产业环节—公司暴露—财务放大}链条的候选主体。下表将链条拆成六类节点；它既用于阅读材料，也用于后续图谱映射与评分输入。节点间关系是研究假设，不等同于收入确认、订单或投资建议。

\begin{table}[htbp]
\centering
\caption{低空产业研究对象的分层与可观察问题}
\label{tab:wty-industry-layers}
\small
\begin{tabularx}{\linewidth}{p{0.18\linewidth}p{0.22\linewidth}X}
\toprule
层次 & 典型对象 & 研究时需要核对的公开问题 \\
\midrule
材料与零部件 & 复合材料、功率器件、传感器 & 是否存在明确技术或供货关系，而非仅有概念关联 \\
动力系统 & 电池、电机、电控、热管理 & 原材料、认证或安全事件能否传导至产品成本与交付 \\
飞行控制与软件 & 飞控、通信、导航、监视 & 事件影响的是设备需求、合规要求还是运行能力 \\
整机与制造 & 无人机、eVTOL、总装与试验 & 取证、产能、交付与应用场景分别处于何种状态 \\
基础设施 & 起降点、充换电、低空智联网 & 政策是否给出建设方向，是否已有可核验实施主体 \\
运营与应用 & 物流、巡检、文旅、应急等 & 场景表述是否只是规划，还是已有公开运行许可与事实 \\
\bottomrule
\end{tabularx}
\end{table}

因此，本计划书把“产业全景”落实为可维护的节点和关系，而非一张静态海报：每一条关系要能回到来源、日期和原文位置；缺乏来源的关系标注为待核验。金融分析只在这条链条已说明的范围内讨论影响强度，不把政策目标直接换算为企业营收或估值。
